% Part: computability
% Chapter: recursive-functions
% Section: non-pr-functions

\documentclass[../../../include/open-logic-section]{subfiles}

\begin{document}

\olfileid{cmp}{rec}{npr}
\olsection{نابنسټيزې بازګشتي تابعې}

بنسټيزې بازګشتي تابعې ټولې هغه تابعې نۀ رانغاړي چې په شهودي ډول محاسبه کېدونکې دي۔ په شهودي ډول بايد څرګنده وي چې د ټولو يوځاييزو بنسټيزو بازګشتي تابعو يو لست~$f_0$، $f_1$، $f_2$،~\dots جوړولے شو، داسې چې د~$y$ پر ننوت د~$f_x$ ارزښت په اغېزمنه توګه محاسبه کړے شو؛ په بله وينا، لاندې تعريف شوې تابع~$g(x,y)$ محاسبه کېدونکې ده:
\[
g(x,y) = f_x(y)
\]
خو بيا لاندې تابع هم محاسبه کېدونکې ده:
\begin{eqnarray*}
h(x) & = & g(x,x) + 1 \\
& = & f_x(x) +1.
\end{eqnarray*}
د هرې بنسټيزې بازګشتي تابع~$f_i$ دپاره د~$h$ او~$f_i$ ارزښتونه په~$i$ کښې سره بېل دي۔ نو~$h$ محاسبه کېدونکې خو بنسټيزه بازګشتي نۀ ده؛ د~$g$ په باب هم همدا ويلے شو۔ دا د کانتور د قطري‌کولو د استدلال يوه «اغېزمنه» بڼه ده۔

د هغو محاسبه کېدونکو تابعو لا څرګندې بېلګې هم ورکولے شو چې بنسټيزې بازګشتي نۀ دي۔ مثلاً نښه~$g^n(x)$ دې~$g(g(\dots g(x)))$ وښيي، چې ټولې~$n$ دانې~$g$ګانې لري؛ او د تابعو لړۍ~$g_0,g_1,\dots$ داسې تعريف کړئ:
\begin{eqnarray*}
g_0(x) & = & x+1 \\
g_{n + 1}(x) & = & g_n^x(x)
\end{eqnarray*}
تاسو ارزولے شئ چې هره تابع~$g_n$ بنسټيزه بازګشتي ده۔ هره ورپسې تابع له مخکينۍ څخه ډېره چټکه لوېږي؛ $g_1(x)$ له~$2x$ سره برابره ده، $g_2(x)$ له~$2^x \cdot x$ سره برابره ده، او~$g_3(x)$ نژدې د~$x$ دانو~$2$ګانو د تواني برج په شان لوېږي۔ د اکرمن--پېټر تابع په اصل کښې همدا تابع~$G(x) = g_x(x)$ ده، او ښودل کېدے شي چې دا له هرې بنسټيزې بازګشتي تابع څخه چټکه لوېږي۔

راځئ د بنسټيزو بازګشتي تابعو د شمېرنې مسئلې ته بېرته راشو۔ په ياد ولرئ چې هرې بنسټيزې بازګشتي تابع ته مو سمبوليکه نښه ټاکلې ده؛ نو د نښو شمېرل بس دي۔ هرې نښې~$F$ ته يو طبيعي عدد~$\#(F)$ په بازګشتي ډول داسې ټاکلے شو:
\begin{eqnarray*}
\#(\Zero) & = & \langle 0 \rangle \\
\#(\Succ) & = & \langle 1 \rangle \\
\#(\Proj{n}{i}) & = & \langle 2, n, i \rangle \\
\#(\fn{Comp}_{k,l}[H,G_0,\dots,G_{k-1}]) & = & \langle 3,k,l,\#(H),\#(G_0),\dots,\#(G_{k-1}) \rangle \\
\#(\fn{Rec}_l[G,H]) & = & \langle 4, l, \#(G), \#(H) \rangle
\end{eqnarray*}
\textbf{د ژباړې د نسخې سمون (OLCMP-015):} سرچينې په دې کوډولو کښې د پخوانۍ برخې ټاکلې صفر او راتلونکې نښې په ناڅاپي ډول په دوو نورو نښو بدلې کړې وې۔ دلته هماغه رسمي نښې بېرته کارول شوې دي؛ د هغو عددي ټاګونه هماغسې پاتې دي۔
دلته دا حقيقت کاروو چې د عددونو هره لړۍ د لړيو د مخکيني کوډ له لارې د يوه طبيعي عدد په توګه ګڼل کېدے شي۔ پايله دا ده چې هرې نښې ته يو طبيعي عددي کوډ ټاکل شوے دے۔ البته ځينې لړۍ (او ځکه ځينې عددونه) له هېڅ نښې سره سمون نۀ لري؛ خو که~$i$ د يوې يوځاييزې بنسټيزې بازګشتي تابع د نښې کوډ وي، $f_i$ هماغه تابع نيسو، او که~$i$ داسې نښه نۀ کوډوي نو د ثابتې صفر تابع ئې نيسو۔ وروستۍ پايله دا ده چې د يوځاييزو بنسټيزو بازګشتي تابعو د شمېرلو يوه څرګنده طريقه لرو۔
\textbf{د ژباړې د نسخې سمون (OLCMP-016):} سرچينې د کوډ ټاکنې پايله په سرچپه لوري بيان کړې وه، حال دا چې پورته بازګشتي معادلې هرې \emph{نښې} ته د لړۍ يو عددي کوډ ټاکي۔ دلته د نقشې لوری او په ورپسې شمېرنه کښې د يوځاييزې نښې شرط څرګند کړل شول۔

(په حقيقت کښې ځينې تابعې، لکه ثابته~$0$ تابع، په لست کښې تر يو ځل زياتې راځي۔ دا يوازې زموږ د کوډولو مصنوعي پايله نۀ ده، بلکې د دې حقيقت پايله هم ده چې ثابته صفر تابع له يوې څخه زياتې نښې لري۔ وروسته به ووينو چې دا تکرارونه په محاسبه کېدونکې توګه نۀ شي مخنيوے کېدے؛ مثلاً هېڅ داسې محاسبه کېدونکې تابع نشته چې پرېکړه وکړي يوه ورکړل شوې نښه ثابته صفر تابع ښيي که نۀ۔)

اوس تابع~$g(x,y)$ د~$f_x(y)$ په توګه اخلو، چې~$f_x$ هماغې شمېرنې ته اشاره کوي چې اوس مو بيان کړه۔ څنګه پوهېږو چې~$g(x,y)$ محاسبه کېدونکې ده؟ په شهودي ډول خبره څرګنده ده: د~$g(x,y)$ د محاسبې دپاره لومړے~$x$ «پرانيزئ» او وګورئ چې ايا دا د يوې يوځاييزې تابع نښه ده که نۀ۔ که وي، پر ننوت~$y$ د هغې تابع ارزښت محاسبه کړئ۔

\begin{tagblock}{TMs}
\begin{digress}
کېدے شي لا مو منلې وي چې (په لږ کار سره!) داسې يو پروګرام، مثلاً په جاوا يا~C++ کښې، ليکلے شو چې دا کار وکړي؛ بيا د چرچ--ټيورينګ اصل ته مراجعه کولے شو، چې وايي هر هغه څه چې په شهودي ډول محاسبه کېدونکې وي، په ټيورينګ ماشين محاسبه کېدے شي۔

البته د~$g(x,y)$ د محاسبه کېدو د ښودلو لا نېغه لار دا ده چې هغه ټيورينګ ماشين په څرګند ډول بيان کړو چې دا تابع محاسبه کوي۔ په ځانګړي ډول، دا به د چرچ--ټيورينګ اصل او شهود ته د مراجعې مخه ونيسي۔ ډېر ژر به دومره وسايل جوړ کړي وي چې د~$g(x,y)$ محاسبه کېدل د محاسبې د داسې مدل په کارولو وښيو چې پر ټيورينګ ماشين \emph{تقليدېدے} شي: يعنې بازګشتي تابعې۔
\end{digress}
\end{tagblock}

\end{document}
